

% Ideas for Introduction:

% Attacking representations in distributed learning
%     MTL happens to be a good case study for this
% Compare with existing attacks in other settings
%     E.g. having each of the prediction heads makes the adversary’s job easier
%     FL where models share the top layer (User-level inference)
% Compare with black-box access to one of the MTL classifiers
% Personalization in distributed learning
%     Embedding is the min info that needs to be shared
%     Multi-task learning is a way to achieve this
%     Threat model with min level of access, black-box to embedding
% Representation/Encoder is the minimal piece of a user’s contribution to the final model
% https://arxiv.org/abs/2309.05505 https://arxiv.org/abs/2002.04758 (both of these are “make FL better” papers not attack papers)
% I.e. If you want to train a model jointly on users’ data, the smallest thing each user must share is the representation updated on their data, which should not be specific to any given user


\section{Introduction}

\emph{Multitask learning} (MTL) has emerged as a powerful paradigm that leverages similarities among multiple learning tasks, each with insufficient samples to train a standalone model, to solve them simultaneously while minimizing data sharing across multiple entities, such as users and organizations.
% The increasing prevalence of decentralized data, scattered across several devices, users, and organizations, has driven the development of methods that %machine learning approaches to 
% effectively train machine learning models over many distinct, yet related, data distributions. 
% % These methods %aim to 
% % enable multiple participants with few samples to jointly train models on their combined data to minimize the sharing of individual data samples and  achieve greater utility through specialization on a given task. 
% %
% Among these methods, \emph{multitask learning (MTL)} has emerged as a powerful paradigm that leverages similarities across multiple tasks %with sparse data 
% to solve them simultaneously. 
MTL accomplishes this goal by learning a \emph{shared representation} that captures common structure between the tasks. Concretely, a shared representation could be a neural network that learns a mapping of the data from all tasks into a shared feature space, where similar data points across tasks cluster together, and task-specific output layers (or heads) operate on these embeddings—that is, the outputs of the shared representation—to make predictions. MTL methods have shown remarkable success in various domains, including computer vision~\cite{girshick2015fastrcnn}, natural language processing~\cite{collobert2008multitasklanguage}, federated learning~\cite{smith2017federated, hanzely2021federatedmixture, mansour2020personalization, ghosh2020clusteredfederated}, drug discovery~\cite{ramsundar2015drugdiscovery}, and financial forecasting~\cite{ghosn1996finance}. For example, MTL can be used to personalize image classification models by learning a shared across many users and locally adapting to each user's small, on-device photo library.

% Talk about how we have to share representations in order to train accurate models that are also personalized


% Justify MTL not by claiming that it offers some flavor of privacy, but rather that it's an effective way to learn models jointly over multiple, sparse tasks
% 1. Why is the threat model (getting representations instead of logits, predictions, etc) realistic/meaningful/justifiable? Is it really white box?
% 2. Does it limit the attacker to get only the representation? 
Despite being designed to capture only generic patterns that can be applied to several downstream tasks, these shared representations can inadvertently leak sensitive information about specific tasks, or underlying data distributions, that they were trained on. The privacy risks that arise are of particular concern when data from several sensitive entities, such as individual users, private organizations, or data silos, are jointly used for training across multiple tasks or non-uniform data distributions. Here, shared representations are often the minimum piece of information that each entity must contribute in order to achieve strong generalization while only contributing a limited number of samples. This point is highlighted in prior work on privacy-preserving collaborative learning~\cite{shen2023sharerep}, where sharing \emph{only} the representation, not task-specific layers, demonstrably improves convergence rates and model performance when learning over multiple parties. 
%  ~\cite{boenisch2023flnotprivate, suri2023subjectmembershipinferenceattacks, nasr2019comprehensive}.


In our work, we study privacy attacks at the level of entire tasks and broadly investigate the privacy risks of jointly learning over multiple tasks by mounting these attacks on the smallest unit of information required to jointly train a model, the shared representation. In particular, we attempt to determine whether, and to what extent, an adversary can infer a given task's inclusion in training given black-box query access to (i.e. the ability to sample embedding vectors from) the shared representation. While there is prior work that explores privacy attacks on MTL, this work is limited to sample-level membership inference and model extraction~\cite{yan2024miamtl} where the adversary can make queries to the task-specific heads of the multitask model and train reference models. 

Towards this goal, we propose a new threat model called \emph{task-inference} and develop black-box attacks with minimal adversarial knowledge to infer the inclusion of a task in MTL. Unlike membership-inference \cite{homer-2008-mi,SankararamanOJH09,dwork-2015-trace,shokri-2016-mi}, task-inference generalizes individual privacy attacks from the sample level to the task level, where the adversary aims to determine the presence of an entire target task, rather than a particular sample, in the training set. Critically, our threat model only assumes that the adversary has access to samples from the target task's distribution. We identify two variants of our threat model, one with a \textbf{strong} adversary who has the specific samples used to train the shared representation and another with a \textbf{weak} adversary who only has independent samples from the target task's distribution. Our experiments demonstrate a notable separation between these two threat models in terms of their attack capabilities,  and we provide analysis of a simplified learning setting to help explain our empirical findings in machine learning. We observe that while both adversaries can mount successful attacks, having access to training samples provides a sizeable advantage. Furthermore, in contrast to prior work \cite{shokri-2016-mi, carlini-2022-lira, 2021-liu-encodermi, chen2023faceauditor, chaudhari-2022-snap, abascal2023tmi, yan2024miamtl}, neither of the task-inference adversaries require auxiliary datasets to train reference, or shadow, models to calibrate their attack. We instead mount attacks by using the key observation that representation models tend to output codependent embedding vectors for samples that come from the same task. 

Because MTL is a general framework that enables learning high utility models by leveraging similarities between related problems, there are several different ways in which tasks can be defined. In this work, we study privacy leakage in two use cases of MTL which we believe are representative of common applications: MTL across many individuals, or users, for \textbf{personalization} and MTL for solving \textbf{multiple related learning problems}. A canonical use of MTL for personalization is photo tagging (i.e. detecting items, attributes, or people in images), where each task corresponds to an individual's mobile device with limited data who cannot train an model on their own. Thus, they learn a shared representation for a common learning problem with other participants in MTL and locally personalize a task-specific layer to their data. In contrast, MTL could also be used to solve several distinct learning problems with limited data but related structure, such as detection of multiple facial attributes in images or multiple topics in text. Here, each task can be seen as a different learning problem with limited data and a unique label space, but each task benefits from learning the structure of other tasks (e.g. facial attributes or topics can be correlated). We observe that the privacy risks that arise from our threat model depend on the specific application of MTL. When using MTL for personalization and tasks are defined as people, inferring inclusion of a task in training can be seen as analogous to membership~\cite{homer-2008-mi} or user-inference~\cite{kandpal2023userinference}. In this setting, the adversary composes the model's outputs over several samples belonging to an individual to determine if any of their data appeared in the dataset. If MTL is used to simultaneously solve multiple learning problems, inferring a task's inclusion poses a similar risk to inferring whether a property~\cite{ateniese-2015-property_inference, hartmann-2023-distributioninference, chaudhari-2022-snap}, or subpopulation that satisfies a common labeling, was present at all in the training dataset. Here, we observe that, similar to property inference, an adversary uses several samples sharing common attributes in the input space (e.g. all text samples discussing the topic "Python") to determine either if, or to what extent, samples satisfying these attributes were included in the training dataset. 

We highlight that the task-inference threat model we propose interpolates those of well-studied privacy attacks on machine learning, and it induces risks comparable to membership, user, and property inference depending on factors like the number of samples per task and the particular definition of tasks. While our analysis in Section~\ref{sec:tracing} offers broad intuition on the efficacy of privacy attacks on groups and non-i.i.d. data, connecting them to sample-level attacks, MTL is a particularly apt regime for studying leakage at the granularity of multiple related samples.

% For example, suppose that MTL is used to learn accurate, \textbf{personalized} classifiers for a common learning problem across sparse data from many different users (e.g. spam detection, autocomplete, etc.). Then, an adversary with the ability to infer whether a task was included in training poses a similar privacy risk to membership-inference~\cite{shokri-2016-mi} and user inference~\cite{kandpal2023userinference} attacks, where inclusion of an individual's contribution to the training dataset is leaked. In contrast, when MTL is used to learn several distinct classifiers with a shared representation to solve \textbf{multiple learning problems} (e.g. facial attribute detection, topic detection, etc.) at once, the privacy risk posed by inferring the inclusion of a task resembles property existence attacks~\cite{chaudhari-2022-snap}, which arise as an extension of the property (or distribution) inference threat model~\cite{ateniese-2015-property_inference}. In the setting where each task is a distinct learning problem, a task-inference adversary can infer whether large subpopulations of the data that satisfy a specific property (or positive labeling) were included during training. 

We comprehensively evaluate our black-box task-inference attacks across vision and language datasets, including CelebA~\cite{liu-2015-celeba} and Stack Overflow~\cite{stackoverflow}, where shared representations are trained using MTL for both personalization across several users and for simultaneously solving multiple learning problems. We also study the factors that lead to task-inference leakage by observing how attack success varies in the model's generalization gap and conducting ablation studies on synthetic data across parameters such as the total number of tasks in the dataset, the embedding dimension, and the number of samples per task. Our findings demonstrate that a purely black-box adversary can successfully infer inclusion of tasks in MTL using only query access to the shared representation, even with \textbf{weak} access to training data samples, reinforcing the generality of our threat model and the separation between \textbf{strong} and \textbf{weak} adversaries observed through our theoretical analysis of tracing attacks.

\subsection*{Our Contributions} We summarize the main contributions of our work as follows:

% Main contributions
% 1. Two new threat models for privacy attacks at the task level. We identify differences between strong and weak adversaries. Mention strong and weak explicitly in the intro. Problem of task inference is novel/New attack on models trained over multiple tasks
% 2. We analyze the threat model on tracing attacks for mean estimation and find that there is a notable separation between the two threat models 
% 3. We propose two purely black-box attacks on ML models trained over multiple tasks (distributions)
% 4. We perform extensive evaluation and [add findings from section 6]
\begin{enumerate}[labelsep=*,leftmargin=15pt]
    \item We formulate novel, black-box threat models for privacy attacks on datasets which are composed of multiple tasks or distributions, focusing on shared representations learned with MTL. Our proposed threat models distinguish between adversaries with \textbf{strong} versus \textbf{weak} access to the shared representation's training data.
    \item We analyze our threat models in the context of tracing attacks for mean estimation over mixtures of Gaussian distributions, and demonstrate a notable separation between the \textbf{strong} and \textbf{weak} adversaries in terms of achievable attack success rates. In addition, we highlight the connections between our tracing attack on tasks and tracing attacks on individual samples. 
    \item We construct two efficient, purely black-box privacy attacks for machine learning models trained over multiple tasks, demonstrating that an adversary can infer the inclusion of a task when they only have \textbf{weak} access to the training data, but can mount significantly more powerful attacks with \textbf{strong} access.
    \item We extensively evaluate both of our attacks on shared representations in machine learning by running experiments in both the vision and language domains for two representative use-cases of MTL, and we empirically study the factors that lead to privacy leakage when learning shared representations over several tasks.
\end{enumerate}
